<artifact artifact_id="abba31cf-352f-4d05-be08-bb3511f20413" artifact_version_id="f1e2d3c4-5a6b-7c8d-9e0f-123456789abc" title="detailed_bit_operations_presentation.tex" contentType="text/latex">
\documentclass{beamer}
% Setting up the document
\usetheme{Madrid}
\usecolortheme{default}
\usepackage[utf8]{inputenc}
\usepackage{CJKutf8}
\usepackage{amsmath}
\usepackage{listings}
\usepackage{xcolor}
% Configuring font for compatibility
\usepackage[T1]{fontenc}
\usepackage{lmodern}
% Configuring listings for code display
\lstset{
basicstyle=\ttfamily\small,
breaklines=true,
frame=single,
numbers=left,
numberstyle=\tiny,
keywordstyle=\color{blue},
commentstyle=\color{gray}
}
% Starting the document
\begin{document}
% Adding title slide
\begin{frame}
\title{位运算及其应用}
\subtitle{深入理解二进制操作与算法技巧}
\author{Generated by Grok}
\date{\today}
\titlepage
\end{frame}
% Adding table of contents
\begin{frame}
\frametitle{目录}
\tableofcontents
\end{frame}
% Section 1: Bit Operations and Properties
\section{位运算及其性质}
\begin{frame}
\frametitle{位运算概述}
\begin{itemize}
\item \textbf{定义}：计算机中数据以二进制存储，位运算是对二进制位直接操作（如加减乘除底层实现）。
\item \textbf{基本运算符}：
\begin{itemize}
\item \texttt{AND (&)}：对应位均为1，结果为1（用于掩码、位清零）。
\item \texttt{OR (|)}：对应位任一为1，结果为1（用于位设置）。
\item \texttt{XOR (^)}：对应位不同，结果为1（用于翻转、唯一元素查找）。
\item \texttt{NOT (~)}：对应位翻转（0↔1）。
\item 左移 (\texttt{<<})：向左移位，相当于乘2（无符号）。
\item 右移 (\texttt{>>})：向右移位，相当于除2。
\end{itemize}
\item \textbf{示例}：35 & 47（二进制：00100011 & 00101111 = 00100011 = 35）。
\end{itemize}
\end{frame}
\begin{frame}
\frametitle{位运算性质}
\begin{itemize}
\item \textbf{与/或性质}：
\begin{itemize}
\item 交换律：$a & b = b & a$，$a | b = b | a$
\item 结合律：$(a & b) & c = a & (b & c)$
\item 分配律：$a & (b | c) = (a & b) | (a & c)$
\item 恒等：$a & 0 = 0$，$a | 0 = a$，$a | \sim a = -1$（全1）
\end{itemize}
\item \textbf{异或性质}（最常用）：
\begin{itemize}
\item $a \oplus a = 0$，$a \oplus 0 = a$
\item 交换/结合：$a \oplus b = b \oplus a$，$(a \oplus b) \oplus c = a \oplus (b \oplus c)$
\item 计数：二进制位中1的个数为奇偶性（用于奇偶校验）。
\end{itemize}
\item \textbf{移位性质}：$a << 1 = 2a$，$a >> 1 = \lfloor a/2 \rfloor$（无符号）。
\end{itemize}
\end{frame}
\begin{frame}
\frametitle{位运算应用示例}
\begin{itemize}
\item \textbf{判断奇偶}：$n & 1$（最低位为1则奇数）。
\item \textbf{交换两数}（无临时变量）：$a = a \oplus b$，$b = a \oplus b$，$a = a \oplus b$。
\item \textbf{代码示例}（C++，交换）：
\begin{lstlisting}[language=C++]
void swap(int& a, int& b) {
a ^= b;
b ^= a;
a ^= b;
}
\end{lstlisting}
\end{itemize}
\end{frame}
% Section 2: Binary Thinking Problems
\section{二进制思维题}
\begin{frame}
\frametitle{二进制思维题概述}
\begin{itemize}
\item \textbf{核心}：利用二进制位表示状态或计数，解决组合/计数问题（如苹果分袋、黄金分割）。
\item \textbf{经典例子1}：苹果分袋（二进制计数）。小明有10个袋子（对应二进制10位），买N（<1024）苹果，用二进制表示N，将苹果放入对应袋子（位为1的袋子放1个）。
\item \textbf{经典例子2}：黄金分割（7天工资）。老板有1段金，每天给1/7。用二进制表示天数：第1天给1/128，第2天给1/64，... 第7天给64/128=1/2。累计正好1/7。
\item \textbf{技巧}：二进制位权（$2^{k-1}$）模拟最小单位计数，避免复杂分数。
\end{itemize}
\end{frame>
\begin{frame}
\frametitle{二进制思维题：LeetCode经典}
\begin{itemize}
\item \textbf{问题}：数组中只出现一次的数字（LeetCode 136）。其他数字出现偶数次。
\item \textbf{思路}：异或性质 $a \oplus a = 0$，$a \oplus 0 = a$。全异或后，结果为唯一数字。
\item \textbf{代码示例}（C++）：
\begin{lstlisting}[language=C++]
int singleNumber(vector<int>& nums) {
int res = 0;
for (int num : nums) res ^= num;  // 时间O(n)，空间O(1)
return res;
}
\end{lstlisting}
\item \textbf{扩展}：牛客“小乐乐改数字”——用位运算将数字转为0/1序列。
\end{itemize}
\end{frame}
% Section 3: Binary Enumeration
\section{二进制枚举}
\begin{frame}
\frametitle{二进制枚举概述}
\begin{itemize}
\item \textbf{定义}：用二进制位（0/1）表示选择状态，枚举$2^n$种子集组合（n≤20，避免超时）。
\item \textbf{应用}：子集求和、组合优化、状态压缩DP预备。
\item \textbf{技巧}：for(int i=0; i<(1<<n); i++)枚举状态；if(i & (1<<j))检查第j位是否选。
\item \textbf{时间复杂度}：O(2^n * n)，适合小n。
\end{itemize}
\end{frame>
\begin{frame}
\frametitle{二进制枚举示例：子集生成}
\begin{itemize}
\item \textbf{问题}：集合{1,2,3}的所有子集。
\item \textbf{枚举过程}：i从0到7（111二进制）。
\begin{itemize}
\item i=0 (000): 空集
\item i=1 (001): {1}
\item i=3 (011): {1,2}
\item i=7 (111): {1,2,3}
\end{itemize}
\item \textbf{代码示例}（C++，输出子集）：
\begin{lstlisting}[language=C++]
vector<int> a = {1,2,3}; int n=3;
for(int i=0; i<(1<<n); i++) {
for(int j=0; j<n; j++) {
if(i & (1<<j)) cout << a[j] << " ";
}
cout << endl;  // 输出所有子集
}
\end{lstlisting}
\end{itemize}
\end{frame}
\begin{frame}
\frametitle{二进制枚举应用：容斥原理}
\begin{itemize}
\item \textbf{问题}：多集合交集计数，用二进制枚举交/并。
\item \textbf{技巧}：位1表示“取交”，位0表示“取并”。公式：$\sum (-1)^{|S|+1} | \cap_{i \in S} A_i |$。
\item \textbf{扩展}：LeetCode子集问题（78），用枚举生成所有组合。
\end{itemize}
\end{frame}
% Section 4: Prefix XOR Sum
\section{前缀异或和}
\begin{frame}
\frametitle{前缀异或和概述}
\begin{itemize}
\item \textbf{定义}：数组a的前缀异或s[i] = a[0] ⊕ a[1] ⊕ ... ⊕ a[i-1]（s[0]=0）。
\item \textbf{核心性质}：区间[l,r]异或 = s[r+1] ⊕ s[l]（异或自逆：x ⊕ x = 0）。
\item \textbf{优势}：预处理O(n)，查询O(1)。用于子数组异或查询、零异或子段计数。
\item \textbf{时间复杂度}：预处理O(n)，多查询O(q)。
\end{itemize}
\end{frame}
\begin{frame}
\frametitle{前缀异或和示例：LeetCode 1310}
\begin{itemize}
\item \textbf{问题}：给定arr和queries[[l,r]]，求每个[l,r]异或值。
\item \textbf{示例}：arr=[1,3,4,8]，queries=[[0,1],[1,2],[0,3]] → [1⊕3=2, 3⊕4=7, 1⊕3⊕4⊕8=14]。
\item \textbf{代码示例}（C++）：
\begin{lstlisting}[language=C++]
vector<int> xorQueries(vector<int>& arr, vector<vector<int>>& queries) {
vector<int> s(arr.size()+1, 0);
for(int i=1; i<=arr.size(); i++) s[i] = s[i-1] ^ arr[i-1];
vector<int> res;
for(auto& q : queries) res.push_back(s[q[1]+1] ^ s[q[0]]);
return res;
}
\end{lstlisting}
\end{itemize}
\end{frame>
\begin{frame}
\frametitle{前缀异或和技巧与扩展}
\begin{itemize}
\item \textbf{技巧}：异或链表（空间优化，双向遍历用前驱⊕后继）。
\item \textbf{应用}：计数异或为0子数组（用map记录前缀出现频次：ans += freq[s[i]]）。
\item \textbf{经典题}：牛客Xorto——枚举优化+前缀异或求最大值。
\item \textbf{注意}：异或无进位，适合位独立计算。
\end{itemize}
\end{frame}
% Ending the document
\end{document}
</artifact>